\documentclass{article}
\usepackage[margin=1in]{geometry}
\usepackage{amsmath,amssymb}
\usepackage[most]{tcolorbox}

\newcommand{\BookName}{Pattern Recognition and Machine Learning}
\newcommand{\BookAuthor}{Christopher M. Bishop}
\newcommand{\ChapterName}{Introduction}
\newcommand{\ExerciseID}{1.4}

\title{Chapter: \ChapterName}
\author{Ahonakpon Guy Gbaguidi}
\date{}

\begin{document}

\maketitle

\begin{center}
  \large\textbf{Exercise \ExerciseID}

  \vspace{0.5em}

  \normalsize\BookName\ - \BookAuthor
\end{center}

\begin{tcolorbox}[breakable,colback=gray!5,colframe=gray!60,title=Solution]
Let's consider a probability density $p_x(x)$ defined over a continuous variable 
$x$ and suppose that we have make a nonlinear change of variable using $x = g(y)$, 
so that the density transforms according  to (1.27):

\begin{align*}
  p_y(y) &= p_x(x) \left| \frac{dx}{dy} \right| \\
  &= p_x(g(y)) \left| g'(y) \right|
\end{align*}

By differentiating (1.27), we have to show that the location $\hat{y}$ of the maximum
of $p_y(y)$ is not in general related to the location $\hat{x}$ of the maximum of $p_x(x)$ 
by the simple functional relation $\hat{x} = g(\hat{y})$ as a consequence of the presence 
of the Jacobian factor $\left| g'(y) \right|$ in (1.27).
\par\medskip

To find the location $\hat{y}$ of the maximum of $p_y(y)$, we need to compute the derivative 
of $p_y(y)$ with respect to $y$ and set it to zero:

\begin{align*}
  \frac{d}{dy} p_y(y) &= 0 \\
  \frac{d}{dy} \left[ p_x(g(y)) \left| g'(y) \right| \right] &= 0
\end{align*}

By applying the product rule of differentiation, we can expand the derivative as follows:

\begin{align*}
  \frac{d}{dy} \left[ p_x(g(y)) \left| g'(y) \right| \right] &= \frac{d}{dy} p_x(g(y)) \cdot \left| g'(y) \right| + p_x(g(y)) \cdot \frac{d}{dy} \left| g'(y) \right| \\
  &= p_x'(g(y)) \cdot g'(y) \cdot \left| g'(y) \right| + p_x(g(y)) \cdot \frac{g''(y)}{\left| g'(y) \right|}
\end{align*}

Setting this equation to zero and solving for $y$ will give us the location of the maximum of $p_y(y)$.
\begin{align*}
  p_x'(g(y)) \cdot g'(y) \cdot \left| g'(y) \right| + p_x(g(y)) \cdot \frac{g''(y)}{\left| g'(y) \right|} &= 0 \\
  p_x'(g(y)) \cdot g'(y) \cdot \left| g'(y) \right| &= - p_x(g(y)) \cdot \frac{g''(y)}{\left| g'(y) \right|} \\
  p_x'(g(y)) \cdot g'(y) \cdot \left| g'(y) \right|^2 &= - p_x(g(y)) \cdot g''(y)
\end{align*}

This shows that the maximum of $p_y(y)$ depends on both the shape of $p_x(x)$ and the transformation $g(y)$. 
So, in general, the point $\hat{y}$ that maximizes $p_y(y)$ is not simply the inverse image of $\hat{x}$,
 because the Jacobian factor $\left| g'(y) \right|$ also changes the density.
\par\medskip

\paragraph{Note from the book:} This shows that the maximum of a probability density (in contrast to a simple 
function) is dependent on the choice of variable.
\par\medskip

Let's verify that, in case of a linear transformation, the location of the maximum transforms 
in the same way as the variable. 
For a linear transformation, we have $g(y) = ay + b$, where $a$ and $b$ are constants.

\par\medskip
The case $a = 0$ is not a valid change of variable, because the transformation is not invertible. 
So we only consider the case $a \neq 0$. In this case, the derivative of $g(y)$ is constant and equal 
to $a$, so we have:

\begin{align*}
  g'(y) = a
\end{align*}

The Jacobian factor is constant and equal to $|a|$. Therefore, the transformed density is given by:

\begin{align*}
  p_y(y) &= p_x(ay + b) \cdot |a|
\end{align*}

To find the location $\hat{y}$ of the maximum of $p_y(y)$, we compute the derivative with respect 
to $y$ and set it to zero:


\begin{align*}
  \frac{d}{dy} p_y(y) &= \frac{d}{dy} \left[ p_x(ay + b) \cdot |a| \right] \\
  &= |a| \cdot p_x'(ay + b) \cdot a \\
  &= a|a| \cdot p_x'(ay + b)
\end{align*}


Setting this equation to zero gives us:

\begin{align*}
  a|a| \cdot p_x'(ay + b) &= 0 \\
  p_x'(ay + b) &= 0 \quad \text{(since $a \neq 0$)}
\end{align*}

Here, we know that $x = g(y) = ay + b$ and $p_x'(x)$ is the derivative of $p_x(x)$ with respect to $x$.

This means that the maximum of $p_y(y)$ occurs when $ay + b = \hat{x}$. Therefore,

\begin{align*}
  \hat{y} = \frac{\hat{x} - b}{a}.
\end{align*}

So, for a linear transformation with $a \neq 0$, the location of the maximum changes in the same way as the variable.
\par\medskip

\paragraph{Personal note:} During the process of solving this exercise, 
I understand that the location of the maximum in maths terms means:

\begin{align*}
  \hat{y} = \arg\max_y p_y(y) \\
  \hat{x} = \arg\max_x p_x(x)
\end{align*}

where $\arg\max$ is the argument of the maximum, which is the value of the variable that maximizes the function.
\end{tcolorbox}

\end{document}
